\chapter{除法与倍数}\label{ch:g5:division}

在\cref{ch:g4:division}中学过的长除法有两项升级：
一是两位数的除数，二是商可以延续到\omterm{def:g4:decimals:point}{小数点}之后——
$3 \div 4$ 终于有了答案 $0.75$。本章以
\omterm{def:g5:division:multiple}{倍数}、\omterm{def:g5:division:multiple}{因数}与最基本的整除速判法作结。

\section{除以两位数}

\begin{method}[两位数除数]\label{met:g5:division:twodigits}
方法与之前相同，只多了一项技能：估计
除数大约能除几次。先列一张除数的小\omterm{def:g5:division:multiple}{倍数}表，
再以 $986 \div 23$ 为例（23 的\omterm{def:g5:division:multiple}{倍数}：23、46、69、92、115、
138、161、184、207）：
\begin{enumerate}
  \item $98$ 个十：$23 \times 4 = 92$ 可以，$23 \times 5 = 115$ 不
        够，商的数字为 $4$，\omterm{def:g3:division:remainder}{余数}为 $6$；
  \item 落下 $6$，得 $66$：$23 \times 2 = 46$ 可以，商的数字为
        $2$，\omterm{def:g3:division:remainder}{余数}为 $20$；
  \item 商为 $42$，\omterm{def:g3:division:remainder}{余数}为 $20$。验算：
        $23 \times 42 + 20 = 966 + 20 = 986$。
\end{enumerate}
\end{method}

\section{小数商}

\begin{example}[不肯停在余数处的除法]\label{ex:g5:division:decimal}
把 $3$ 个披萨分给 $4$ 个孩子：$3 \div 4$ 的商是 $0$，
\omterm{def:g3:division:remainder}{余数}是 $3$——毫无用处！把除法继续到\omterm{def:g4:decimals:point}{小数点}之后：
\begin{enumerate}
  \item $3$ 个一 $= 30$ 个十分之一；$30 \div 4$ 得 $7$ 个十分之一，余
        $2$ 个十分之一；
  \item $2$ 个十分之一 $= 20$ 个百分之一；$20 \div 4 = 5$ 个百分之一，
        \omterm{def:g3:division:remainder}{余数}为 $0$。
\end{enumerate}
于是 $3 \div 4 = 0.75$，每个孩子分得 $0.75$ 个披萨——
也就是四分之三，正如\omterm{def:g4:fractions:def}{分数}已经说明的那样
（见\cref{prop:g5:fractions:decimal}）。
\end{example}

\begin{method}[把除法延续到小数点之后]\label{met:g5:division:continue}
当个位已经除完、仍有\omterm{def:g3:division:remainder}{余数}时：
\begin{enumerate}
  \item 在商中写上\omterm{def:g4:decimals:point}{小数点}；
  \item 在\omterm{def:g3:division:remainder}{余数}后面添一个零（个位变成十分位，十分位
        变成百分位\,\dots），然后继续除；
  \item \omterm{def:g3:division:remainder}{余数}为 $0$ 时停止——或数字开始
        无限循环，如 $1 \div 3 = 0.333\dots$，则给出
        \omterm{def:g4:numbers:round}{四舍五入}后的近似值。
\end{enumerate}
\end{method}

\begin{example}\label{ex:g5:division:continue}
$27 \div 6$：商为 $4$，\omterm{def:g3:division:remainder}{余数}为 $3$；写上\omterm{def:g4:decimals:point}{小数点}，$30$ 个十分之一
$\div\ 6 = 5$，于是 $27 \div 6 = 4.5$。\quad
$22 \div 8$：商为 $2$，\omterm{def:g3:division:remainder}{余数}为 $6$；再 $60 \div 8 = 7$ 余 $4$；
再 $40 \div 8 = 5$，于是 $22 \div 8 = 2.75$。
\end{example}

\section{倍数与因数}

\begin{definition}[倍数、因数]\label{def:g5:division:multiple}
$6$ 的\emph{倍数}\index{倍数}，就是把它的乘法表
一直延续下去所得的数：$0, 6, 12, 18, 24, \dots$。当 $24$ 是 $6$ 的倍数时，
也可以说 $6$ 是 $24$ 的\emph{因数}\index{因数}，
或说 $24$ \emph{能被 $6$ 整除}：
除法 $24 \div 6$ 没有\omterm{def:g3:division:remainder}{余数}。
\end{definition}

\begin{proposition}[整除速判法]\label{prop:g5:division:criteria}
不做除法，即可判断一个数是否能被下列各数整除：
\begin{itemize}
  \item 能被 $2$ 整除：末位数字是 $0$、$2$、$4$、$6$ 或 $8$ 时
        （即\omterm{def:g3:numbers:evenodd}{偶数}）；
  \item 能被 $5$ 整除：末位数字是 $0$ 或 $5$；
  \item 能被 $10$ 整除：末位数字是 $0$；
  \item 能被 $3$ 整除：它的\emph{各位数字之和}能被
        $3$ 整除时（例如\ $741$：$7 + 4 + 1 = 12$，能被整除——事实上
        $741 = 3 \times 247$）。
\end{itemize}
\end{proposition}
\admitted

\begin{omfigure}
\begin{tikzpicture}[scale=0.5]
  \foreach \n in {0,...,29} {
    \pgfmathtruncatemacro{\r}{int(\n/10)}
    \pgfmathtruncatemacro{\c}{mod(\n,10)}
    \pgfmathtruncatemacro{\lbl}{\n+1}
    \pgfmathtruncatemacro{\mtwo}{mod(\lbl,2)}
    \pgfmathtruncatemacro{\mthree}{mod(\lbl,3)}
    \ifnum\mtwo=0
      \fill[omDef!25] (\c,-\r) ++(-0.45,-0.45) rectangle ++(0.9,0.9);
    \fi
    \ifnum\mthree=0
      \draw[omThm, thick] (\c,-\r) circle (0.42);
    \fi
    \node[font=\footnotesize] at (\c,-\r) {\lbl};
  }
  \node[font=\small, below] at (4.5,-2.8)
    {蓝色方块：$2$ 的倍数；\ \ 红色圆圈：
    $3$ 的倍数；\ \ 两者兼有：$6$ 的倍数};
\end{tikzpicture}

{\small $1$ 到 $30$ 这些数中，$2$ 和 $3$ 的\omterm{def:g5:division:multiple}{倍数}形成了一些图案
——而同时带有两种标记的数（$6$、$12$、$18$、
$24$、$30$）恰好是 $6$ 的\omterm{def:g5:division:multiple}{倍数}。}
\end{omfigure}

\begin{example}\label{ex:g5:division:criteria}
$534$ 能被 $2$ 整除吗？末位是 $4$，能。能被 $5$ 整除吗？不能。能被 $3$ 整除吗？
$5 + 3 + 4 = 12$，能。因此 $534$ 也能被 $6$ 整除（既能被 $2$
\emph{又}能被 $3$ 整除）——验算：$534 = 6 \times 89$。
\end{example}

\section{练习}

\begin{exercise}[$\star$]\label{exo:g5:division:1}
用竖式计算下列除法（先列出\omterm{def:g5:division:multiple}{倍数}表）：
$851 \div 23$; \quad $704 \div 32$.
\end{exercise}

\begin{exercise}[$\star$]\label{exo:g5:division:2}
计算下列商（写成准确的小数）：$9 \div 2$；\quad $27 \div 4$；
\quad $33 \div 5$; \quad $21 \div 8$.
\end{exercise}

\begin{exercise}[$\star$]\label{exo:g5:division:3}
四个朋友平分一笔 $54$ 的餐厅账单，每人
应付多少？（继续除到\omterm{def:g4:decimals:point}{小数点}以后。）
\end{exercise}

\begin{exercise}[$\star$]\label{exo:g5:division:4}
计算 $10 \div 3$，商继续除到\omterm{def:g4:decimals:point}{小数点}后三位，
你观察到了什么？请写出精确到百分位的近似值。
\end{exercise}

\begin{exercise}[$\star$]\label{exo:g5:division:5}
列出下列各数：$7$ 的\omterm{def:g5:division:multiple}{倍数}（$70$ 以内）；$18$ 的\omterm{def:g5:division:multiple}{因数}；
$24$ 的\omterm{def:g5:division:multiple}{因数}。
\end{exercise}

\begin{exercise}[$\star$]\label{exo:g5:division:6}
能被 $2$ 整除吗？能被 $5$ 整除吗？能被 $10$ 整除吗？能被 $3$ 整除吗？分别用对应的判别法检验下列各数：
$470$; \quad $735$; \quad $8\,001$; \quad $1\,314$.
\end{exercise}

\begin{exercise}[$\star$]\label{exo:g5:division:7}
求出 $60$ 与 $80$ 之间所有
既能被 $3$ \emph{又}能被 $5$ 整除的数。（这两个条件合在一起，
相当于能被哪一个数整除？）
\end{exercise}

\begin{exercise}[$\star$]\label{exo:g5:division:8}
一根长 $7.2$ m 的彩带剪成 $6$ 段相等的小段，
每段长多少？（小数除以整数：先分整米，
再分十分之一。）
\end{exercise}

\begin{exercise}[$\star$]\label{exo:g5:division:9}
$1\,000$ 颗弹珠装袋，每袋装 $24$ 颗。能装满多少袋，
还剩多少颗？（除数是两位数。）
\end{exercise}

\begin{exercise}[$\star\star$]\label{exo:g5:division:10}
$5$ 个朋友平分 $7$ 条巧克力，能否恰好分完？
每人分得的量写成小数。另外，$7$ 个朋友平分
$5$ 条巧克力，情况有什么不同？
\end{exercise}

\begin{exercise}[$\star\star$]\label{exo:g5:division:11}
利用数字和的判别法，求最小的数字 $d$，使
$52d$（末位是 $d$ 的三位数）能被 $3$ 整除。然后
求出所有满足条件的数字 $d$。
\end{exercise}
